% !TEX program = pdflatex
% !BIB program = biber
% =========================================================================
%  Praca dyplomowa – Wydział Inżynierii Środowiska i Energetyki PK
%
%  Kompilacja:  latexmk            (ustawienia w pliku latexmkrc)
%  lub ręcznie: pdflatex praca → biber praca → pdflatex praca → pdflatex praca
%
%  Struktura katalogów:
%    strony/         – strona tytułowa (tu wpisuje się dane pracy: autor,
%                      temat, promotor...), oświadczenie, wykaz oznaczeń
%    rozdzialy/      – rozdziały pracy (jeden plik = jeden rozdział)
%    zalaczniki/     – załączniki
%    rysunki/        – pliki graficzne; rysunki/logo/ – oficjalne logo PK i WIŚiE
%    dane/, kod/     – pliki CSV i kody programów wczytywane do pracy
%    literatura.bib  – baza opisów bibliograficznych
%    wisiepk.cls     – ustawienia formatowania (zwykle bez zmian)
% =========================================================================
\documentclass[jednostronnie]{wisiepk}
% Dostępne opcje klasy:
%   dwustronnie – wydruk dwustronny (marginesy lustrzane, puste strony)
%   arial       – krój bezszeryfowy zamiast Times
%   bibpodzial  – literatura podzielona na publikacje / normy / akty prawne / WWW

\addbibresource{literatura.bib}

% Miejsce na własne pakiety i polecenia.

\begin{document}

% =========================================================================
%  STRONA TYTUŁOWA
%
%  1. Uzupełnij dane w bloku poniżej – trafią one na stronę tytułową,
%     do oświadczenia i do metadanych PDF.
%  2. Część „Układ strony” nie wymaga edycji.
% =========================================================================

% ------------------------- DANE PRACY ------------------------------------
\autor{Imię Nazwisko}
\plecautora{M}              % M – formy męskie, K – żeńskie (w oświadczeniu)

\tytulPL{Temat pracy dyplomowej w~języku polskim}

\rodzajpracy{inżynierska}   % inżynierska / magisterska
\formastudiow{stacjonarne}  % stacjonarne / niestacjonarne
\stopienstudiow{I}          % I – studia inżynierskie, II – magisterskie
\kierunek{Inżynieria Środowiska}  % np. Energetyka, Geoinformatyka
\specjalnosc{brak}

\promotor{dr inż. Imię Nazwisko}
\recenzent{dr hab. inż. Imię Nazwisko}  % puste – pole „Recenzent” znika

\miejsce{Kraków}
\rok{2026}

\zgodaASAP{TAK}             % zgoda na udostępnienie pracy w ASAP: TAK / NIE
% -------------------------------------------------------------------------

\metadanepdf

% ------------------------- UKŁAD STRONY ----------------------------------
\newgeometry{top=1.7cm, bottom=1.3cm, left=2.2cm, right=2.2cm}
\thispagestyle{empty}
\begingroup
\setstretch{1}
\fontfamily{phv}\selectfont         % krój bezszeryfowy, jak na wzorze wydziału
\setlength{\parindent}{0pt}
\newcommand{\etykieta}[1]{{\fontsize{9.5}{11}\selectfont\color{black!55}#1}}

% Nagłówek: logo wydziału i sygnet PK
\logowydzialu\hfill\logopk\par
\vspace{5mm}
{\color{wisieniebieski}\rule{\linewidth}{0.6pt}}\par

\vspace*{\stretch{1.3}}
\begin{center}
  {\fontsize{22}{27}\selectfont\color{pkgranat}\dana{autor}\par}
  \vspace{9mm}
  {\fontsize{17}{22}\selectfont\bfseries\color{pkgranat}\dana{tytulPL}\par}
  \vspace{6mm}
  {\fontsize{11}{13}\selectfont praca \dana{rodzajpracy}\par}

  \vspace{13mm}
  {\fontsize{11}{17}\selectfont
    \textbf{studia \dana{formastudiow} \dana{stopienstudiow} stopnia}\\
    \etykieta{kierunek studiów:} \dana{kierunek}\\
    \etykieta{specjalność:} \dana{specjalnosc}\par}
\end{center}

\vspace*{\stretch{1}}
\begin{flushright}
  \etykieta{Promotor:}\\[2pt]
  {\fontsize{11.5}{14}\selectfont\bfseries\dana{promotor}}
  \jeslidana{recenzent}{\\[9mm]
    \etykieta{Recenzent:}\\[2pt]
    {\fontsize{11.5}{14}\selectfont\dana{recenzent}}}
\end{flushright}

\vspace*{\stretch{1.1}}
\begin{center}
  {\fontsize{11.5}{14}\selectfont\dana{miejsce}, \dana{rok}}
\end{center}
\vspace*{\stretch{0.5}}

% Stopka z danymi wydziału
{\color{wisieniebieski}\rule{\linewidth}{0.4pt}}\par
\vspace{2.5mm}
\begin{minipage}[b]{0.62\linewidth}
  \fontsize{7.5}{9.5}\selectfont
  {\bfseries\color{pkgranat}Wydział Inżynierii Środowiska i~Energetyki}\\
  ul. Warszawska 24, 31-155 Kraków\\
  tel. 12 628 28 01, s-0@pk.edu.pl
\end{minipage}\hfill
\begin{minipage}[b]{0.36\linewidth}
  \raggedleft\fontsize{11}{13}\selectfont\bfseries\color{wisieniebieski}%
  www.wisie.pk.edu.pl
\end{minipage}\par
\endgroup
\restoregeometry
\cleardoublepage
      % ← tu uzupełnij dane pracy
% =========================================================================
%  Oświadczenie autora pracy.
%  Treść odpowiada formularzowi stosowanemu na PK. Przed wydrukiem porównaj
%  ją z aktualnym wzorem z dziekanatu – przepisy i numery Dz.U. zmieniają się
%  po każdym ogłoszeniu tekstu jednolitego ustawy.
%  Formy męskie/żeńskie oraz wybór TAK/NIE ustawia się w strony/tytulowa.tex.
% =========================================================================
\thispagestyle{empty}
\begingroup
\small\setstretch{1.05}
\setlist[enumerate]{leftmargin=*, itemsep=3pt}

\begin{center}
  {\large\bfseries OŚWIADCZENIE O SAMODZIELNYM WYKONANIU\\ PRACY DYPLOMOWEJ}
\end{center}
\vspace{4pt}

Oświadczam, że przedkładana przeze mnie praca dyplomowa \dana{rodzajpracy}
została napisana przeze mnie samodzielnie. Jednocześnie oświadczam, że
ww.~praca:
\begin{enumerate}
  \item nie narusza praw autorskich w~rozumieniu ustawy z~dnia 4~lutego 1994~r.
    o~prawie autorskim i~prawach pokrewnych (Dz.U. z~2025~r. poz.~24) oraz
    dóbr osobistych chronionych prawem cywilnym, a~także nie zawiera danych
    i~informacji, które \formaMK{uzyskałem}{uzyskałam}{uzyskałem/am}
    w~sposób niedozwolony,
  \item nie była wcześniej podstawą żadnej innej procedury związanej
    z~nadawaniem tytułów zawodowych, stopni lub tytułów naukowych.
\end{enumerate}

Jednocześnie wyrażam zgodę na:
\begin{enumerate}
  \item poddanie mojej pracy kontroli za pomocą systemu Antyplagiat oraz na
    umieszczenie tekstu pracy w~bazie danych uczelni, w~celu ochrony go przed
    nieuprawnionym wykorzystaniem. Oświadczam, że
    \formaMK{zostałem poinformowany}{zostałam poinformowana}{zostałem/am poinformowany/a}
    i~wyrażam zgodę, by system Antyplagiat porównywał tekst mojej pracy
    z~tekstem innych prac znajdujących się w~bazie danych uczelni,
    z~tekstami dostępnymi w~zasobach światowego Internetu oraz z~bazą
    porównawczą systemu Antyplagiat,
  \item to, aby moja praca pozostała w~bazie danych uczelni przez okres
    wynikający z~przepisów prawa. Oświadczam, że
    \formaMK{zostałem poinformowany}{zostałam poinformowana}{zostałem/am poinformowany/a}
    i~wyrażam zgodę, że tekst mojej pracy stanie się elementem porównawczej
    bazy danych uczelni, która będzie wykorzystywana, w~tym także udostępniana
    innym podmiotom, na zasadach określonych przez uczelnię, w~celu
    dokonywania kontroli antyplagiatowej prac dyplomowych/doktorskich,
    a~także innych tekstów, które powstaną w~przyszłości.
    \miejscenapodpis
  \item Wyrażam zgodę na udostępnianie mojej pracy dyplomowej w~Akademickim
    Systemie Archiwizacji Prac na PK do celów naukowo-badawczych
    z~poszanowaniem przepisów ustawy o~prawie autorskim i~prawach pokrewnych
    (Dz.U. z~2025~r. poz.~24).
    \par\vspace{4pt}\hfill\makebox[6cm]{\wyborTakNie}
    \miejscenapodpis
\end{enumerate}

\vspace{4pt}
\textit{Jednocześnie przyjmuję do wiadomości, że w~przypadku stwierdzenia
popełnienia przeze mnie czynu polegającego na przypisaniu sobie autorstwa
istotnego fragmentu lub innych elementów cudzej pracy, lub ustalenia
naukowego, Rektor PK stwierdzi nieważność postępowania w~sprawie nadania mi
tytułu zawodowego (art.~77 ust.~5 ustawy z~dnia 20~lipca 2018~r. Prawo
o~szkolnictwie wyższym i~nauce (Dz.U. z~2024~r. poz.~1571, z~późn.~zm.)).}
\miejscenapodpis

\dopisekSkreslenia
\endgroup
\cleardoublepage


\tableofcontents
% =========================================================================
%  Wykaz ważniejszych oznaczeń (opcjonalny). Jeśli praca go nie wymaga,
%  usuń % =========================================================================
%  Wykaz ważniejszych oznaczeń (opcjonalny). Jeśli praca go nie wymaga,
%  usuń % =========================================================================
%  Wykaz ważniejszych oznaczeń (opcjonalny). Jeśli praca go nie wymaga,
%  usuń \input{strony/oznaczenia} z pliku praca.tex.
%  Symbole wpisuje się bez znaków $ – kolumna jest w trybie matematycznym.
% =========================================================================
\begin{oznaczenia}
  \grupaoznaczen{Litery łacińskie}
  \oznaczenie{c_p}{ciepło właściwe przy stałym ciśnieniu}{\unit{J/(kg.K)}}
  \oznaczenie{d}{grubość warstwy przegrody}{\unit{m}}
  \oznaczenie{D}{średnica wewnętrzna przewodu}{\unit{m}}
  \oznaczenie{L}{długość odcinka przewodu}{\unit{m}}
  \oznaczenie{\Delta p}{strata ciśnienia}{\unit{Pa}}
  \oznaczenie{R}{opór cieplny}{\unit{m^2.K/W}}
  \oznaczenie{U}{współczynnik przenikania ciepła}{\unit{W/(m^2.K)}}
  \oznaczenie{w}{średnia prędkość przepływu}{\unit{m/s}}

  \grupaoznaczen{Litery greckie}
  \oznaczenie{\lambda}{współczynnik przewodzenia ciepła}{\unit{W/(m.K)}}
  \oznaczenie{\lambda_\mathrm{f}}{współczynnik oporów liniowych (tarcia)}{--}
  \oznaczenie{\rho}{gęstość}{\unit{kg/m^3}}
  \oznaczenie{\theta_\mathrm{e}}{temperatura powietrza zewnętrznego}{\unit{\degreeCelsius}}

  \grupaoznaczen{Skróty}
  \oznaczenie{\text{c.o.}}{centralne ogrzewanie}{}
  \oznaczenie{\text{c.w.u.}}{ciepła woda użytkowa}{}
  \oznaczenie{\text{OZE}}{odnawialne źródła energii}{}
\end{oznaczenia}
 z pliku praca.tex.
%  Symbole wpisuje się bez znaków $ – kolumna jest w trybie matematycznym.
% =========================================================================
\begin{oznaczenia}
  \grupaoznaczen{Litery łacińskie}
  \oznaczenie{c_p}{ciepło właściwe przy stałym ciśnieniu}{\unit{J/(kg.K)}}
  \oznaczenie{d}{grubość warstwy przegrody}{\unit{m}}
  \oznaczenie{D}{średnica wewnętrzna przewodu}{\unit{m}}
  \oznaczenie{L}{długość odcinka przewodu}{\unit{m}}
  \oznaczenie{\Delta p}{strata ciśnienia}{\unit{Pa}}
  \oznaczenie{R}{opór cieplny}{\unit{m^2.K/W}}
  \oznaczenie{U}{współczynnik przenikania ciepła}{\unit{W/(m^2.K)}}
  \oznaczenie{w}{średnia prędkość przepływu}{\unit{m/s}}

  \grupaoznaczen{Litery greckie}
  \oznaczenie{\lambda}{współczynnik przewodzenia ciepła}{\unit{W/(m.K)}}
  \oznaczenie{\lambda_\mathrm{f}}{współczynnik oporów liniowych (tarcia)}{--}
  \oznaczenie{\rho}{gęstość}{\unit{kg/m^3}}
  \oznaczenie{\theta_\mathrm{e}}{temperatura powietrza zewnętrznego}{\unit{\degreeCelsius}}

  \grupaoznaczen{Skróty}
  \oznaczenie{\text{c.o.}}{centralne ogrzewanie}{}
  \oznaczenie{\text{c.w.u.}}{ciepła woda użytkowa}{}
  \oznaczenie{\text{OZE}}{odnawialne źródła energii}{}
\end{oznaczenia}
 z pliku praca.tex.
%  Symbole wpisuje się bez znaków $ – kolumna jest w trybie matematycznym.
% =========================================================================
\begin{oznaczenia}
  \grupaoznaczen{Litery łacińskie}
  \oznaczenie{c_p}{ciepło właściwe przy stałym ciśnieniu}{\unit{J/(kg.K)}}
  \oznaczenie{d}{grubość warstwy przegrody}{\unit{m}}
  \oznaczenie{D}{średnica wewnętrzna przewodu}{\unit{m}}
  \oznaczenie{L}{długość odcinka przewodu}{\unit{m}}
  \oznaczenie{\Delta p}{strata ciśnienia}{\unit{Pa}}
  \oznaczenie{R}{opór cieplny}{\unit{m^2.K/W}}
  \oznaczenie{U}{współczynnik przenikania ciepła}{\unit{W/(m^2.K)}}
  \oznaczenie{w}{średnia prędkość przepływu}{\unit{m/s}}

  \grupaoznaczen{Litery greckie}
  \oznaczenie{\lambda}{współczynnik przewodzenia ciepła}{\unit{W/(m.K)}}
  \oznaczenie{\lambda_\mathrm{f}}{współczynnik oporów liniowych (tarcia)}{--}
  \oznaczenie{\rho}{gęstość}{\unit{kg/m^3}}
  \oznaczenie{\theta_\mathrm{e}}{temperatura powietrza zewnętrznego}{\unit{\degreeCelsius}}

  \grupaoznaczen{Skróty}
  \oznaczenie{\text{c.o.}}{centralne ogrzewanie}{}
  \oznaczenie{\text{c.w.u.}}{ciepła woda użytkowa}{}
  \oznaczenie{\text{OZE}}{odnawialne źródła energii}{}
\end{oznaczenia}


\chapter{Cel i zakres pracy}
\label{chap:cel}

% Rozdział powinien zmieścić się na jednej stronie.

Rozdział otwiera się jednoznacznym sformułowaniem celu pracy, czyli tego, co
w~jej ramach zostało zaprojektowane, obliczone, zbadane lub ocenione.
Przykładowo: \emph{Celem pracy jest opracowanie koncepcji instalacji
grzewczej budynku jednorodzinnego zasilanej z~powietrznej pompy ciepła oraz
porównanie tego rozwiązania z~kotłem gazowym pod względem kosztów
eksploatacji i~emisji CO\textsubscript{2}.}

Następnie wymienia się zadania, których realizacja doprowadziła do
osiągnięcia celu. Zakres pracy obejmował:
\begin{itemize}
  \item przegląd literatury oraz przepisów i~norm dotyczących
        projektowanego systemu,
  \item obliczenie projektowego obciążenia cieplnego budynku zgodnie
        z~normą PN-EN 12831-1 \cite{pn12831},
  \item dobór urządzeń oraz obliczenia hydrauliczne instalacji,
  \item analizę ekonomiczną i~środowiskową rozpatrywanych wariantów,
  \item sformułowanie wniosków i~zaleceń projektowych.
\end{itemize}

W~tej części warto zapowiedzieć układ pracy, np.: \emph{W~rozdziale drugim
omówiono stan wiedzy na temat \dots, rozdział trzeci zawiera obliczenia
\dots, a~w~rozdziale czwartym przedstawiono \dots} Należy tu również
wymienić programy komputerowe wykorzystane do obliczeń, symulacji
i~przygotowania rysunków, np.\ QGIS \cite{qgis} albo EPANET \cite{epanet22}.

\chapter{Wprowadzenie}
\label{chap:wprowadzenie}

Wprowadzenie osadza temat pracy w~szerszym kontekście: wyjaśnia, dlaczego
podjęty problem jest istotny z~punktu widzenia inżynierii środowiska lub
energetyki, jakie są obecne tendencje i~wymagania prawne oraz jakie
rozwiązania są stosowane w~praktyce. Przykładowo, w~pracy dotyczącej
ciepłownictwa można przywołać koncepcję sieci czwartej generacji
\cite{lund2014} oraz przegląd stanu ciepłownictwa na świecie
\cite{werner2017}, a~w~pracy o~budynkach -- wymagania dyrektywy
w~sprawie charakterystyki energetycznej budynków \cite{epbd2024}.

Poniżej zebrano wymagania edytorskie, które formatka spełnia automatycznie.
Przy pisaniu pracy nie trzeba ich ustawiać ręcznie -- wystarczy korzystać
z~poleceń opisanych w~rozdziale~\ref{chap:redakcja}.

\section{Wymagania edytorskie}
\label{sec:wymagania}

Podstawowe parametry składu zestawiono w~tabeli~\ref{tab:parametry}.
Każdy rozdział główny rozpoczyna się od nowej strony. Strona tytułowa jest
wliczana do numeracji, lecz numer nie jest na niej drukowany.

\begin{table}[htbp]
  \centering
  \caption{Parametry formatowania pracy dyplomowej}
  \label{tab:parametry}
  \begin{tabularx}{\linewidth}{@{}>{\raggedright\arraybackslash}p{4.4cm}X@{}}
    \toprule
    \textbf{Element} & \textbf{Ustawienie} \\
    \midrule
    Format strony      & A4, druk jedno- lub dwustronny \\
    Marginesy          & lewy (wewnętrzny) \qty{3.5}{cm}, pozostałe \qty{2.5}{cm} \\
    Krój i stopień     & Krój typu Times lub Arial, \qty{11}{pt} \\
    Interlinia         & 1,25 wiersza (tekst), pojedyncza (podpisy, tabele, kod) \\
    Akapit             & wyjustowany, wcięcie pierwszego wiersza \qty{0.63}{cm} \\
    Tytuły rozdziałów  & \qty{14}{pt}, pogrubione, numeracja „1.” \\
    Podrozdziały       & \qty{13}{pt} i~\qty{12}{pt}, pogrubione, numeracja „1.1.” i~„1.1.1.” \\
    Podpisy            & kursywa, rysunek -- pod obiektem, tabela -- nad obiektem \\
    Numeracja obiektów & w~obrębie rozdziału: rys.~3.1, tabela~3.1, wzór~(3.1) \\
    Numer strony       & na dole strony, wyśrodkowany \\
    \bottomrule
  \end{tabularx}
\end{table}

\section{Druk i oprawa}

Jeżeli jeden z~egzemplarzy ma być wydrukowany dwustronnie (np.\ w~miękkiej
oprawie), należy w~pliku \texttt{praca.tex} zmienić opcję klasy
z~\texttt{jednostronnie} na \texttt{dwustronnie}. Marginesy staną się
wówczas lustrzane (szerszy margines zawsze od strony grzbietu), rozdziały
będą zaczynały się na stronach nieparzystych, a~wstawione strony puste nie
otrzymają numeru. Wersję do archiwizacji w~systemie ASAP zwykle składa się
jednostronnie -- warto potwierdzić to z~promotorem.

\chapter{Zasady redagowania tekstu}
\label{chap:redakcja}

Rozdział pokazuje, jak w~tej formatce wstawia się wzory, rysunki, tabele
i~fragmenty kodu. Kod źródłowy rozdziału (plik
\texttt{rozdzialy/03\_redakcja.tex}) można traktować jako ściągę --
wystarczy skopiować potrzebny fragment i~podmienić treść.

\section{Styl wypowiedzi}
\label{sec:styl}

Pracę pisze się w~formie bezosobowej. Unika się pierwszej osoby liczby
pojedynczej i~mnogiej, a~także zwrotów potocznych. Porównanie wariantów
poprawnych i~niepoprawnych zawiera tabela~\ref{tab:styl}.

\begin{table}[htbp]
  \centering
  \caption{Przykłady formy bezosobowej}
  \label{tab:styl}
  \begin{tabularx}{\linewidth}{@{}Xc@{}}
    \toprule
    \textbf{Zdanie} & \textbf{Ocena} \\
    \midrule
    Wyniki obliczeń strat ciśnienia zestawiono w~tabeli~3.4. & poprawnie \\
    Wyniki obliczeń strat ciśnienia przedstawia tabela~3.4.  & poprawnie \\
    Wyniki obliczeń strat ciśnienia zestawiłem w~tabeli~3.4. & niepoprawnie \\
    W~naszej pracy policzyliśmy straty ciśnienia.            & niepoprawnie \\
    \bottomrule
  \end{tabularx}
\end{table}

Kilka zasad typograficznych, o~które warto zadbać od początku:
\begin{itemize}
  \item jednoliterowe spójniki i~przyimki (a, i, o, u, w, z) nie mogą
        zostawać na końcu wiersza -- po nich wpisuje się tyldę, np.\
        \verb|w~budynku|,
  \item liczbę od jednostki oddziela się spacją niełamliwą; najprościej
        użyć polecenia \verb|\qty|, np.\ \verb|\qty{2.5}{kW}| daje
        \qty{2.5}{kW} (z~przecinkiem dziesiętnym),
  \item zakresy liczbowe zapisuje się półpauzą: \verb|10--15| daje 10--15,
        a~w~połączeniu z~jednostką: \qtyrange{10}{15}{\degreeCelsius},
  \item cudzysłów polski ma postać „tekst” (\verb|,,tekst''|),
  \item skróty wyjaśnia się przy pierwszym użyciu, np.\ odnawialne źródła
        energii (OZE), i~zbiera w~wykazie oznaczeń.
\end{itemize}

\section{Wzory}
\label{sec:wzory}

Wzory wyrównuje się do środka, a~ich numery -- w~nawiasach okrągłych -- do
prawego marginesu. Numeracja jest prowadzona w~obrębie rozdziału. Wzór
stanowi część zdania, dlatego po nim stawia się przecinek lub kropkę.
Wielkości występujące we wzorze objaśnia się bezpośrednio pod nim
(środowisko \texttt{gdzie}), podając jednostki w~nawiasach kwadratowych.

\subsection{Wzór pojedynczy}

Współczynnik przenikania ciepła przegrody złożonej z~warstw jednorodnych
oblicza się zgodnie z~normą PN-EN ISO 6946 \cite{pn6946} ze wzoru
\begin{equation}
  U = \frac{1}{R_\mathrm{si} + \sum\limits_{i=1}^{n}\dfrac{d_i}{\lambda_i}
      + R_\mathrm{se}},
  \label{eq:U}
\end{equation}
\begin{gdzie}
  \zm{R_\mathrm{si},\ R_\mathrm{se}}{opór przejmowania ciepła na
      powierzchni wewnętrznej i~zewnętrznej}{\unit{m^2.K/W}}
  \zm{d_i}{grubość $i$-tej warstwy}{\unit{m}}
  \zm{\lambda_i}{obliczeniowy współczynnik przewodzenia ciepła materiału
      $i$-tej warstwy}{\unit{W/(m.K)}}
  \zm{n}{liczba warstw przegrody}{}
\end{gdzie}

Do wzoru odwołuje się poleceniem \verb|\eqref|, np.: \emph{wartość $U$
obliczona ze wzoru~\eqref{eq:U} dla ściany z~listingu~\ref{lst:U} wynosi
ok.\ \qty{0.23}{W/(m^2.K)}}. Zmienne zapisuje się kursywą, a~indeksy
opisowe (np.\ \emph{si} -- \emph{surface internal}) -- pismem prostym:
\verb|R_\mathrm{si}|. Dzięki pakietowi \texttt{icomma} przecinek
dziesiętny we wzorach nie tworzy zbędnego odstępu: $\theta_\mathrm{e} =
-20$~°C, $\eta = 0,92$.

\subsection{Układ kilku wzorów}

Kilka powiązanych zależności wyrównuje się względem znaku równości
(środowisko \texttt{align}). Spadek ciśnienia w~przewodzie opisuje wzór
Darcy'ego--Weisbacha~\eqref{eq:dw}, a~współczynnik oporów liniowych dla
przepływu burzliwego -- równanie Colebrooka--White'a~\eqref{eq:cw}:
\begin{align}
  \Delta p &= \lambda_\mathrm{f}\,\frac{L}{D}\,\frac{\rho\,w^2}{2},
  \label{eq:dw}\\
  \frac{1}{\sqrt{\lambda_\mathrm{f}}} &= -2\log_{10}\left(
    \frac{2,51}{\mathit{Re}\,\sqrt{\lambda_\mathrm{f}}}
    + \frac{k}{3,71\,D}\right),
  \label{eq:cw}
\end{align}
\begin{gdzie}
  \zm{\Delta p}{strata ciśnienia na długości przewodu}{\unit{Pa}}
  \zm{\lambda_\mathrm{f}}{współczynnik oporów liniowych}{}
  \zm{L,\ D}{długość i~średnica wewnętrzna przewodu}{\unit{m}}
  \zm{\rho}{gęstość wody}{\unit{kg/m^3}}
  \zm{w}{średnia prędkość przepływu}{\unit{m/s}}
  \zm{\mathit{Re}}{liczba Reynoldsa}{}
  \zm{k}{chropowatość bezwzględna ścianki}{\unit{m}}
\end{gdzie}

\section{Rysunki}
\label{sec:rysunki}

Rysunkiem jest każda ilustracja: wykres, schemat, mapa, zdjęcie czy zrzut
ekranu -- w~podpisie wszystkie oznacza się skrótem „Rys.”. Rysunek
wyśrodkowuje się, a~podpis (kursywą, bez kropki na końcu) umieszcza pod
nim. \textbf{Każdy rysunek musi zostać omówiony w~tekście}, a~odwołanie
powinno poprzedzać sam rysunek. Ilustracje rastrowe powinny mieć
rozdzielczość co najmniej \qty{300}{dpi}; wykresy najlepiej wstawiać jako
grafikę wektorową (PDF) albo rysować bezpośrednio w~\LaTeX-u. Jeżeli
rysunek pochodzi z~literatury lub powstał na podstawie cudzych danych,
w~podpisie podaje się źródło.

\subsection{Wykres rysowany w~\LaTeX-u}

Przykładową krzywą grzewczą, czyli zależność temperatury wody
zasilającej i~powrotnej od temperatury powietrza zewnętrznego, pokazano na
rys.~\ref{fig:krzywa}. Wykres powstał w~pakiecie \texttt{pgfplots}, więc
czcionki i~przecinki dziesiętne są spójne z~tekstem pracy.

\begin{figure}[htbp]
  \centering
  \begin{tikzpicture}
    \begin{axis}[
      width=0.82\linewidth, height=6.2cm,
      xmin=-20, xmax=20, ymin=15, ymax=75,
      x dir=reverse,
      xlabel={Temperatura zewnętrzna $\theta_\mathrm{e}$ [°C]},
      ylabel={Temperatura wody [°C]},
      grid=major, grid style={black!15},
      legend pos=north west, legend cell align=left,
      every axis plot/.append style={thick, samples=60, domain=-20:19},
    ]
      \addplot[black]        {20 + 50*((20 - x)/40)^0.77};
      \addplot[black, dashed]{20 + 35*((20 - x)/40)^0.77};
      \legend{zasilanie, powrót}
    \end{axis}
  \end{tikzpicture}
  \caption{Przykładowa krzywa grzewcza dla parametrów obliczeniowych
    70/55\,°C (dane umowne)}
  \label{fig:krzywa}
\end{figure}

\subsection{Schemat blokowy}

Proste schematy można rysować w~pakiecie TikZ. Rysunek~\ref{fig:schemat}
przedstawia uproszczony schemat systemu ciepłowniczego.

\begin{figure}[htbp]
  \centering
  \begin{tikzpicture}[
      node distance=8mm and 9mm,
      blok/.style={draw, thick, rounded corners=2pt, align=center,
                   minimum width=2.5cm, minimum height=1.2cm, font=\small},
      zas/.style={-{Latex[length=2.2mm]}, thick},
      pow/.style={-{Latex[length=2.2mm]}, thick, dashed}]
    \node[blok] (zrodlo) {Źródło\\ciepła};
    \node[blok, right=of zrodlo] (siec) {Sieć\\ciepłownicza};
    \node[blok, right=of siec]   (wezel) {Węzeł\\cieplny};
    \node[blok, right=of wezel]  (odb) {Instalacje\\c.o. i~c.w.u.};
    \draw[zas] ([yshift=3mm]zrodlo.east) -- ([yshift=3mm]siec.west);
    \draw[zas] ([yshift=3mm]siec.east)   -- ([yshift=3mm]wezel.west);
    \draw[zas] ([yshift=3mm]wezel.east)  -- ([yshift=3mm]odb.west);
    \draw[pow] ([yshift=-3mm]odb.west)   -- ([yshift=-3mm]wezel.east);
    \draw[pow] ([yshift=-3mm]wezel.west) -- ([yshift=-3mm]siec.east);
    \draw[pow] ([yshift=-3mm]siec.west)  -- ([yshift=-3mm]zrodlo.east);
  \end{tikzpicture}
  \caption{Schemat systemu ciepłowniczego (linia ciągła -- zasilanie,
    przerywana -- powrót)}
  \label{fig:schemat}
\end{figure}

\subsection{Rysunek z~pliku}

Do wstawiania plików graficznych służy polecenie \verb|\includegraphics|
(rys.~\ref{fig:plik}), a~same pliki umieszcza się w~katalogu
\texttt{rysunki/}. Rozszerzenie pliku można pominąć. Użyte tu obrazy zastępcze (pakiet \texttt{mwe}) są
dostępne w~każdej dystrybucji \LaTeX-a -- należy je zastąpić własnymi
ilustracjami.

\begin{figure}[htbp]
  \centering
  \includegraphics[width=0.55\linewidth]{example-image-a}
  \caption{Rysunek wstawiony z~pliku graficznego (źródło: opracowanie
    własne na podstawie \cite{pvgis})}
  \label{fig:plik}
\end{figure}

\subsection{Rysunek złożony}

Kilka powiązanych ilustracji można zestawić obok siebie jako jeden rysunek
z~podpisami cząstkowymi (rys.~\ref{fig:zlozony}). Do części rysunku
odwołuje się tak samo: rys.~\ref{fig:zlozony-a}.

\begin{figure}[htbp]
  \centering
  \begin{subfigure}[b]{0.45\linewidth}
    \centering
    \includegraphics[width=\linewidth]{example-image-b}
    \caption{stan przed modernizacją}
    \label{fig:zlozony-a}
  \end{subfigure}\hfill
  \begin{subfigure}[b]{0.45\linewidth}
    \centering
    \includegraphics[width=\linewidth]{example-image-c}
    \caption{stan po modernizacji}
    \label{fig:zlozony-b}
  \end{subfigure}
  \caption{Przykład rysunku złożonego z~dwóch części}
  \label{fig:zlozony}
\end{figure}

\section{Tabele}
\label{sec:tabele}

Tabelę poprzedza podpis umieszczony \emph{nad} nią. Tak jak rysunek, musi
ona zostać przywołana i~skomentowana w~tekście. Jednostki podaje się
w~nagłówkach kolumn, a~liczby w~kolumnie wyrównuje względem przecinka
dziesiętnego (kolumna typu \texttt{S} z~pakietu \texttt{siunitx}).

\subsection{Tabela wpisana ręcznie}

Orientacyjne wartości współczynnika przewodzenia ciepła wybranych
materiałów budowlanych zestawiono w~tabeli~\ref{tab:lambda}.

\begin{table}[htbp]
  \centering
  \caption{Współczynnik przewodzenia ciepła wybranych materiałów (wartości
    orientacyjne)}
  \label{tab:lambda}
  \begin{tabular}{@{}l S[table-format=4.0] S[table-format=1.3]@{}}
    \toprule
    \textbf{Materiał} & {\textbf{Gęstość} $\rho$ [\unit{kg/m^3}]}
                      & {$\lambda$ [\unit{W/(m.K)}]} \\
    \midrule
    beton zwykły                          & 2400 & 1.70  \\
    mur z~cegły ceramicznej pełnej        & 1800 & 0.77  \\
    tynk cementowo-wapienny               & 1850 & 0.82  \\
    drewno sosnowe (w~poprzek włókien)    & 550  & 0.16  \\
    wełna mineralna                       & 40   & 0.035 \\
    styropian EPS                         & 15   & 0.040 \\
    \bottomrule
  \end{tabular}
\end{table}

\subsection{Tabela wczytana z~pliku CSV}

Dane z~arkusza kalkulacyjnego nie muszą być przepisywane ręcznie. Wystarczy
zapisać je jako plik CSV (w~polskiej wersji Excela separatorem jest średnik,
a~znakiem dziesiętnym przecinek) i~wczytać pakietem \texttt{pgfplotstable}.
Tabela~\ref{tab:csv} powstała w~ten sposób z~pliku
\texttt{dane/zuzycie\_ciepla.csv}.

\begin{table}[htbp]
  \centering
  \caption{Miesięczne zużycie ciepła w~budynku przykładowym (dane umowne)}
  \label{tab:csv}
  \pgfplotstabletypeset[
    col sep=semicolon,
    /pgf/number format/read comma as period,
    columns/miesiac/.style={string type, column name={Miesiąc},
                            column type={c}},
    columns/t_zewn/.style={column name={$\theta_\mathrm{e}$ [°C]},
                           fixed, fixed zerofill, precision=1,
                           column type={r}},
    columns/Q_co/.style={column name={$Q_\text{c.o.}$ [GJ]},
                         fixed, fixed zerofill, precision=1,
                         column type={r}},
    columns/Q_cwu/.style={column name={$Q_\text{c.w.u.}$ [GJ]},
                          fixed, fixed zerofill, precision=1,
                          column type={r}},
    every head row/.style={before row=\toprule, after row=\midrule},
    every last row/.style={after row=\bottomrule},
  ]{dane/zuzycie_ciepla.csv}
\end{table}

Z~tego samego pliku można od razu sporządzić wykres
(rys.~\ref{fig:csv}) -- po zmianie danych w~pliku CSV zarówno tabela, jak
i~wykres zaktualizują się przy kolejnej kompilacji.

\pgfplotstableread[col sep=semicolon]{dane/zuzycie_ciepla.csv}\tabCieplo
\begin{figure}[htbp]
  \centering
  \pgfplotsset{/pgf/number format/read comma as period}
  \begin{tikzpicture}
    \begin{axis}[
      width=0.9\linewidth, height=6cm,
      ybar stacked, bar width=11pt, ymin=0,
      xtick=data, xticklabels from table={\tabCieplo}{miesiac},
      xlabel={Miesiąc}, ylabel={Zużycie ciepła [GJ]},
      ymajorgrids, grid style={black!15},
      legend style={at={(0.5,0.97)}, anchor=north, legend columns=2,
                    draw=none, font=\small},
      enlarge x limits=0.05,
    ]
      \addplot[fill=black!55, draw=black]
        table[x expr=\coordindex, y=Q_co]{\tabCieplo};
      \addplot[fill=black!15, draw=black]
        table[x expr=\coordindex, y=Q_cwu]{\tabCieplo};
      \legend{c.o., c.w.u.}
    \end{axis}
  \end{tikzpicture}
  \caption{Wykres sporządzony na podstawie danych z~tabeli~\ref{tab:csv}}
  \label{fig:csv}
\end{figure}

\section{Fragmenty kodu}
\label{sec:kod}

Krótkie fragmenty programów wstawia się w~środowisku \texttt{lstlisting}
(listing~\ref{lst:U}), dłuższe -- z~pliku poleceniem
\verb|\lstinputlisting| i~najlepiej w~załączniku (listing~\ref{lst:dp}
w~załączniku~\ref{zal:kod}). Polskie znaki w~komentarzach są obsługiwane.

\begin{lstlisting}[language=Python, caption={Obliczanie współczynnika
  przenikania ciepła przegrody}, label={lst:U}]
def wspolczynnik_U(warstwy, R_si=0.13, R_se=0.04):
    """Zwraca współczynnik przenikania ciepła U [W/(m2K)].

    warstwy - lista krotek (grubość d [m], lambda [W/(mK)])
    """
    R_t = R_si + R_se + sum(d / lam for d, lam in warstwy)
    return 1.0 / R_t

# Ściana: tynk 1,5 cm, cegła pełna 25 cm, EPS 15 cm, tynk 0,5 cm
sciana = [(0.015, 0.82), (0.25, 0.77), (0.15, 0.040), (0.005, 0.82)]
print(f"U = {wspolczynnik_U(sciana):.3f} W/(m2K)")
\end{lstlisting}

\section{Wyliczenia i~przypisy}

Wyliczenia punktowane stosuje się dla elementów równorzędnych, numerowane
-- gdy liczy się kolejność, np.\ dla etapów obliczeń:
\begin{enumerate}
  \item zebranie danych o~budynku i~lokalizacji,
  \item obliczenie strat ciepła przez przenikanie i~wentylację:
    \begin{enumerate}
      \item dla poszczególnych pomieszczeń,
      \item dla całego budynku,
    \end{enumerate}
  \item dobór źródła ciepła.
\end{enumerate}
Przypisy dolne służą do krótkich uzupełnień, które przerywałyby tok
wywodu\footnote{Przypisów nie należy używać zamiast wykazu literatury --
źródła podaje się w~nawiasach kwadratowych w~tekście.}.

\chapter{Korzystanie ze źródeł}
\label{chap:zrodla}

Każda informacja, wzór, rysunek, tabela lub dane zaczerpnięte z~cudzego
opracowania muszą zostać opatrzone odwołaniem do źródła. Opisy
bibliograficzne przechowuje się w~pliku \texttt{literatura.bib}, a~wykaz
literatury na końcu pracy tworzy się automatycznie -- trafiają do niego
wyłącznie pozycje przywołane w~tekście, w~kolejności pierwszego cytowania.

\section{Sposób cytowania}
\label{sec:cytowanie}

Odwołanie ma postać numeru w~nawiasie kwadratowym (polecenie
\verb|\cite{klucz}|). Umieszcza się je \textbf{przed} kropką kończącą
zdanie:
\begin{itemize}
  \item \emph{Pompy ciepła mogą w~znacznym stopniu ograniczyć zużycie gazu
        ziemnego do ogrzewania budynków \cite{iea2022}.} -- poprawnie,
  \item \emph{Pompy ciepła mogą w~znacznym stopniu ograniczyć zużycie gazu
        ziemnego do ogrzewania budynków. \cite{iea2022}} -- niepoprawnie.
\end{itemize}

Inne często potrzebne warianty:
\begin{itemize}
  \item kilka źródeł naraz -- \verb|\cite{lund2014,werner2017,iea2022}|
        daje \cite{lund2014,werner2017,iea2022}; kolejne numery są
        automatycznie łączone w~zakres,
  \item wskazanie strony lub rozdziału -- \verb|\cite[s.~25]{duffie2013}|
        daje \cite[s.~25]{duffie2013},
  \item nazwisko autora w~zdaniu -- \verb|\textcite{werner2017}|:
        \emph{\textcite{werner2017} przedstawił przegląd \dots};
        przy więcej niż dwóch autorach pojawia się skrót „i~in.”:
        \emph{\textcite{lund2014} zdefiniowali \dots}
\end{itemize}

Jeżeli promotor wymaga stylu autor--rok, np.\ \emph{(Werner, 2017)},
w~pliku \texttt{wisiepk.cls} należy zmienić \texttt{style=numeric-comp} na
\texttt{style=authoryear} i~używać polecenia \verb|\parencite|.

\section{Przykłady opisów bibliograficznych}
\label{sec:przyklady-bib}

Tabela~\ref{tab:typy-bib} zestawia najczęściej cytowane rodzaje źródeł
wraz z~typem wpisu w~pliku \texttt{literatura.bib} i~przykładem
z~wykazu literatury tej formatki.

\begin{table}[htbp]
  \centering
  \caption{Rodzaje źródeł i~odpowiadające im typy wpisów}
  \label{tab:typy-bib}
  \begin{tabularx}{\linewidth}{@{}>{\raggedright\arraybackslash}X
                                   >{\ttfamily}l c@{}}
    \toprule
    \textbf{Rodzaj źródła} & \textbf{\normalfont\bfseries Typ wpisu} & \textbf{Przykład} \\
    \midrule
    książka (podręcznik, monografia)       & @book          & \cite{pelech2009,duffie2013,metcalf2014} \\
    rozdział w~pracy zbiorowej             & @incollection  & \cite{nuorkivi2016} \\
    artykuł w~czasopiśmie                  & @article       & \cite{lund2014,werner2017} \\
    referat konferencyjny                  & @inproceedings & \cite{crawley2005} \\
    raport, instrukcja, dokumentacja       & @report        & \cite{iea2022,epanet22} \\
    norma krajowa lub międzynarodowa       & @standard      & \cite{pn12831,pn6946,pn14001,iso52000} \\
    ustawa, rozporządzenie, dyrektywa      & @legislation   & \cite{ustawaPOS,rozpWT,epbd2024} \\
    strona internetowa, serwis             & @online        & \cite{pvgis,tlm,isap} \\
    zbiór danych                           & @dataset       & \cite{imgw} \\
    oprogramowanie                         & @software      & \cite{qgis} \\
    karta katalogowa producenta            & @manual        & \cite{katalog} \\
    praca dyplomowa                        & @thesis        & \cite{praca2024} \\
    \bottomrule
  \end{tabularx}
\end{table}

\subsection{Normy}

W~opisie normy podaje się jej pełne oznaczenie \emph{wraz z~datą
wydania} (np.\ PN-EN 12831-1:2017-08), tytuł, wydawcę i~rok. Cytuje się
wersję faktycznie wykorzystaną w~obliczeniach, a~jej aktualność warto
sprawdzić w~katalogu Polskiego Komitetu Normalizacyjnego. Przykładowy
wpis pokazano na listingu~\ref{lst:bib-norma}.

\begin{lstlisting}[language={}, caption={Wpis normy w~pliku
  literatura.bib}, label={lst:bib-norma}]
@standard{pn12831,
  number       = {PN-EN 12831-1:2017-08},
  title        = {Charakterystyka energetyczna budynków -- Metoda
                  obliczania projektowego obciążenia cieplnego --
                  Część 1: Projektowe obciążenie cieplne, Moduł M3-3},
  organization = {Polski Komitet Normalizacyjny},
  location     = {Warszawa},
  year         = {2017},
}
\end{lstlisting}

\subsection{Akty prawne}

Ustawy i~rozporządzenia cytuje się według \emph{aktualnego tekstu
jednolitego} ogłoszonego w~Dzienniku Ustaw, z~dopiskiem „z~późn.~zm.”, jeśli
po jego ogłoszeniu wprowadzono zmiany. Stan prawny i~numer tekstu
jednolitego najłatwiej sprawdzić w~Internetowym Systemie Aktów Prawnych
\cite{isap}. Akty prawa Unii Europejskiej cytuje się według Dziennika
Urzędowego UE \cite{epbd2024}.

\subsection{Źródła internetowe}

Przy stronach internetowych obowiązkowo podaje się adres URL oraz datę
dostępu (pole \texttt{urldate}), ponieważ treść strony może się zmienić
(listing~\ref{lst:bib-www}). Należy korzystać przede wszystkim ze
źródeł instytucjonalnych (urzędy, instytuty badawcze, organizacje
międzynarodowe). Encyklopedie redagowane przez użytkowników, fora i~blogi
nie są wiarygodnym źródłem w~pracy dyplomowej.

\begin{lstlisting}[language={}, caption={Wpis strony internetowej
  w~pliku literatura.bib}, label={lst:bib-www}]
@online{pvgis,
  author       = {{European Commission, Joint Research Centre}},
  title        = {Photovoltaic Geographical Information System (PVGIS)},
  url          = {https://joint-research-centre.ec.europa.eu/...},
  urldate      = {2026-10-09},
}
\end{lstlisting}

Nazwę instytucji występującej w~roli autora ujmuje się w~podwójne
nawiasy klamrowe, aby nie została potraktowana jak imię i~nazwisko.
Datę dostępu zapisuje się w~formacie RRRR-MM-DD -- w~wykazie literatury
zostanie wydrukowana w~formie polskiej.

\chapter{Podsumowanie i~wnioski}
\label{chap:wnioski}

W~podsumowaniu w~zwięzły sposób przypomina się cel pracy i~ocenia, czy oraz
w~jakim stopniu został on osiągnięty. Wnioski powinny wynikać bezpośrednio
z~przeprowadzonych obliczeń, analiz lub badań -- nie wprowadza się tu
nowych informacji ani odwołań do literatury, których wcześniej nie
omówiono.

Wnioski wygodnie jest sformułować w~punktach, np.:
\begin{enumerate}
  \item Projektowe obciążenie cieplne budynku wynosi \qty{8.4}{kW}, co
        pozwala na zastosowanie pompy ciepła o~mocy nominalnej \dots
  \item Roczny koszt eksploatacji wariantu z~pompą ciepła jest o~\dots\%
        niższy niż wariantu z~kotłem gazowym.
  \item Największy wpływ na wynik analizy ekonomicznej ma \dots
\end{enumerate}

Jeżeli praca ma charakter projektowy, należy podać zalecenia wynikające
z~obliczeń. Na końcu można wskazać kierunki dalszych prac, np.\
rozszerzenie analizy o~symulację całoroczną albo pomiary na obiekcie
rzeczywistym.


\wykazliteratury
\spisrysunkow
\spistabel

\zalaczniki
\chapter{Kod programu do obliczania strat ciśnienia}
\label{zal:kod}

Załączniki oznacza się kolejnymi wielkimi literami (A, B, C, \dots).
Umieszcza się w~nich materiały uzupełniające: obszerne wyniki obliczeń,
rysunki techniczne, karty katalogowe czy kod programów. Program
z~listingu~\ref{lst:dp} realizuje obliczenia według
wzorów~\eqref{eq:dw} i~\eqref{eq:cw}.

\lstinputlisting[language=Python, caption={Obliczanie strat ciśnienia
  w~przewodzie (plik kod/straty\_cisnienia.py)},
  label={lst:dp}]{kod/straty_cisnienia.py}


\end{document}
